% ===== 原文 PDF pp.24--27（印刷 pp.23--26）=====
% 译者：tr-02；tag 区间 2.38--2.41；图 3

% [原文 p.23]
$T(\rho)$ 是规范群 $G$ 的表示 $\rho$ 中的二次 Casimir 不变量（Dynkin 指数），其归一化使得
$T(\mathrm{adj})$ 等于 $G$ 的对偶 Coxeter 数。\footnote{对 $G=\mathrm{U}(1)$ 的情形，有
$T=\frac{1}{2}q^2$，其中 $q$ 是在该规范 $\mathrm{U}(1)$ 群下的荷。}

来自 4d $\cN=2$ 超对称的非重整化性质使规范耦合的跑动得到极大的简化：它断言 beta 函数是
1 圈精确的，因此我们可以放心地忽略高圈修正 $\mathcal O(g^5)$。这是全纯性的一个简单推论。
具体地说，在某一特定的重整化方案下，超势是手征超场的全纯函数；这些超场中包括背景手征超场，
其 VEV 被看作背景复化耦合 $\tau$。于是，由于微扰重整化不能依赖 $\theta_{\mathrm{YM}}$
（它与一个拓扑项相联系），$n$ 圈阶的这类贡献必须携带一个因子 $\Im(\tau)^{1-n}$；而除非
$n=1$，这个因子并不全纯。因此我们得出结论：规范耦合 $g$ 的重整化在微扰论中是 1 圈精确的，
\footnote{1 圈 beta 函数在全纯方案下对 $\cN=1$ 与 $\cN=2$ SCFT 都精确；但在前一种情形下，
对 Kähler 势的修正会重整化物理耦合（即在 NSVZ 方案中）[124--127]。}
相应地 $\tau$ 取如下形式
\begin{equation}
  \tau(E)=\tau_{\mathrm{UV}}-\frac{b}{2\pi i}\log\left(\frac{E}{\Lambda_{\mathrm{UV}}}\right)+\dots\,,
  \tag{2.38}
\end{equation}
其中只有非微扰贡献没有被指定。

刻画非平凡 RG 流的一个重要量是\textbf{动力学标度} $\Lambda$，其定义为
\begin{equation}
  \Lambda^b=E^b e^{2\pi i\tau(E)}\,,
  \tag{2.39}
\end{equation}
它在微扰意义下是 RG 不变的，即 $\partial\Lambda/\partial E=0$。它也被称为\textbf{嬗变标度}，
即一圈耦合发散的那个标度；因此更高圈与非微扰效应都必须被考虑在内。等价地说，$\Lambda$
标志着规范理论强耦合相与弱耦合相之间的转变。

利用动力学标度 $\Lambda$（可以把它看作某个背景手征超场的标量分量），我们可以写下与 (2.38)
的一圈跑动以及全纯性都相容的、关于耦合 $\tau(E)$ 的最一般表达式，
\begin{equation}
  \tau(E)=-\frac{b}{2\pi i}\log\left(\frac{E}{\Lambda}\right)+\sum_{n>0}a_n\left(\frac{\Lambda}{E}\right)^{bn}\,,
  \tag{2.40}
\end{equation}
其中对 $n>0$ 的求和捕捉来自 $n$ 个瞬子的非微扰贡献，即

% [原文 p.24]
\[
  e^{-S_{\mathrm{inst}}}\sim\left(\frac{\Lambda}{E}\right)^{b}.
\]
对于 $\cN=2$ 超对称的场内容，$b$ 化简为如下组合
\begin{equation}
  b=2\big(T(\mathrm{adj})-T(\rho_{\text{matter}})\big)\,,
  \tag{2.41}
\end{equation}
其中 $\rho_{\text{matter}}$ 统称超多重态 $H_{\cN=2}$ 所处的 $G$ 表示。另一方面，瞬子贡献的
系数 $a_n$ 对一个一般理论来说是出了名地难以计算。然而，正如我们将在第 3 节看到的，
Seiberg--Witten 理论 [20,21] 提供了一条优雅的途径，只需少量物理输入就能确定这些非微扰效应。

% [Tutor Note: (2.39) 的 RG 不变性可直接由 (2.38) 验证：把 (2.38) 写成 $2\pi i\tau(E)=2\pi i\tau_{\mathrm{UV}}-b\log(E/\Lambda_{\mathrm{UV}})$，两边取指数后乘 $E^{b}$，得 $E^{b}e^{2\pi i\tau(E)}=\Lambda_{\mathrm{UV}}^{b}e^{2\pi i\tau_{\mathrm{UV}}}$，与 $E$ 无关。这正是原文所说的「微扰意义下 RG 不变」——(2.38) 的 $\dots$ 项（非微扰修正）会破坏这一不变性，故 (2.40) 要把它们单独写成 $\sum_{n>0}a_n(\Lambda/E)^{bn}$。]
% [Tutor Note: $\Im(\tau)^{1-n}$ 这个论证的关键在于「全纯」的定义：$\tau$ 与 $\bar\tau$ 是全纯坐标与反全纯坐标，任何只含 $\Im\tau$ 的非平凡幂次的因子都同时依赖 $\tau$ 与 $\bar\tau$，故不全纯。$n=1$ 时该因子为 $\Im(\tau)^0=1$，唯一允许的情形。这解释了为什么高圈修正（$n\ge2$）在全纯方案下不出现，而不需要逐圈计算。]

\begin{exercisebox}{习题 2.3}
  经典的 $\mathrm{U}(1)_r$ 对称性在一般的 $\cN=2$ 规范理论中是反常的。请找出
  $\mathrm{U}(1)_r$ 反常以及无反常的剩余对称性。
\end{exercisebox}

\begin{exercisebox}{习题 2.4}
  论证：在 (2.40) 中乘在 $a_n$ 上的那些项，是对规范耦合 $\tau$ 的非微扰修正中唯一相关的项。
\end{exercisebox}

% [Tutor Note: 习题 2.3 的答案梗概（供自检）：$\mathrm{U}(1)_r$ 的反常来自 ABJ 型三角形图，其反常系数为 $2b$（$b$ 即 (2.41) 的一圈系数）；因此该对称性被破缺为 $\Z_{2b}$。这个结论原文在 p.28 的脚注 18 中明确给出，可对照阅读。习题 2.4 则是在问：为什么 (2.40) 中 $(\Lambda/E)^{bn}$ 这类项（即 $e^{-nS_{\mathrm{inst}}}$）是仅有的非微扰修正——理由是全纯性加单值性只允许 $\Lambda^{b n}$ 这样的整幂次出现，其它函数形式会在 $\tau\to\tau+1$ 下破坏良好定义。]

\SUB{2.5}{$\cN=2$ SCFT 的拉格朗日量}

在介绍了 4d $\cN=2$ 拉格朗日量与 RG 流的基础之后，人们也许会问：它们如何能帮助我们研究
一般的 $\cN=2$ SCFT？毕竟我们已经知道，（已知的）4d $\cN=2$ SCFT 大多数是强耦合的，
没有直接的拉格朗日量描述。尽管如此，4d $\cN=2$ 拉格朗日量在从紫外或红外视角研究
4d $\cN=2$ SCFT 时都能派上用场。

从 IR 视角看，可以预期某些 4d $\cN=2$ 拉格朗日量会作为带有 $\cN=2$ 保持变形的 SCFT 的
有效描述出现。这些变形可能来自打开相关算子，或移到超对称的真空模空间上。它们一般给出
弱耦合的 $\cN=2$ EFT，并将在第 3 节详细描述。

从 UV 视角看，从 $\cN=2$ 拉格朗日量出发、通向 SCFT 的情形有两种：

\para{$b=0$：}此时规范理论拉格朗日量是共形不变的，从而定义一个 SCFT。特别地，规范耦合 $g$
是恰当边缘的，并参数化该 SCFT 的共形流形。一些简单的例子包括 $\cN=4$ SYM 以及

% [原文 p.25]
图 3a 中所示的 $\cN=2$ 共形 SQCD。

由于 $b$ 由所关心理论的物质内容决定，我们可以通过调整物质内容使 $b=0$，以此对共形
拉格朗日量进行分类。特别地，聚焦于特殊酉规范群 $\otimes_i\SU(N_i)$ 的情形，并假定物质内容
由双基础表示的超多重态给出，则共形箭图拉格朗日量有一个完整的分类，它与 Dynkin 图及仿射
Dynkin 图的分类相吻合。例如，$E_6$ 仿射 Dynkin 图及相应的共形箭图拉格朗日量示于图 3b，
其中每个圆形节点承载一个规范群，每条边住着一个双基础表示的超多重态。最一般的规范群与
物质内容的情形已在 [128] 中分类。

共形拉格朗日量的一个明显优点是：它们能直接促成一系列 SCFT 可观测量的具体计算，例如用局域化
计算超对称球面配分函数——这些配分函数对复化规范耦合有非平凡的依赖 [76--78,129]。

\para{$b>0$：}$\cN=2$ 拉格朗日量在 IR 中导向有趣 SCFT 的另一种情形，是 $\cN=2$ 理论
渐近自由。在这类情形中，超共形对称性并不显式，而只是在 IR 中涌现。人们找到这类涌现 SCFT
的办法，是调节拉格朗日量中某些保持 $\cN=2$ 的参数，例如可能的质量参数、某些受保护算子的
VEV，甚至动力学标度 $\Lambda$ 本身。作为一个例子（示于图 3c），$\SU(2)$ SYM 理论耦合一个
基础表示的超多重态，以及纯 $\SU(3)$ SYM 理论，在 CB 上都拥有一个由 $(A_1,A_2)$
Argyres--Douglas（AD）理论 [22] 描述的超共形不动点。\footnote{我们将在第 4 节（Class
$\cS$）或第 5 节（IIB 型）中讨论如何以自顶向下的方式实现这些理论。}

把这类渐近自由理论与涌现的 IR SCFT 联系起来，乍看像是猜谜游戏，但确有系统性的工具，
它们利用渐近自由拉格朗日量的 IR EFT。其策略是在 CB 上寻找具有涌现标度不变性的区域
（这些区域对应着调节保持 $\cN=2$ 的参数）。最初的 AD 理论正是这样被发现的 [22,23]。

从这些渐近自由的 UV 拉格朗日量中提取 SCFT 可观测量是可能的，但要困难得多。除了调节为在
IR 中定位 SCFT 点所需的参数之外，还必须小心处理由流动产生的退耦部门。在具体理论中已取得
一些进展，例如可参见 [130--132]。

% [原文 p.26]
\begin{center}
  \includegraphics[width=0.82\linewidth]{fig3.png}
\end{center}
\figtitle{3}{显式（上面两幅）与涌现（下面一幅）SCFT 的例子。}

% [Tutor Note: 图 3 内的英文标注是原图固有的，未作改动。读法：圆形节点承载规范群，方形节点表示味对称群（与原文「每个圆形节点承载一个规范群、每条边住着一个双基础表示超多重态」的约定一致）；(b) 中的数字 $n,2n,3n$ 是各节点上的 $\SU(N)$ 的 $N$，构成 $E_6$ 仿射 Dynkin 图（中心节点度数 3）。(c) 的两支箭头表示两个「不同」的 UV 理论——$\SU(2)$ 耦合一个基础超多重态、以及纯 $\SU(3)$ SYM——在调参后于 CB 上流到同一个 IR 超共形不动点，即 $(A_1,A_2)$ AD 理论；这正是原文「emergent SCFT」的含义。]
% [Tutor Note: 注意 $b=0$ 与 $b>0$ 两种情形的地位完全不同：$b=0$ 时共形对称性在 UV 就显式存在（拉格朗日量本身共形不变，$g$ 是恰当边缘耦合）；$b>0$ 时 UV 理论渐近自由、只有标度不变性在 IR 涌现，必须调参才能落到不动点上。原文用「manifest (top two) or emergent (bottom one)」概括图 3 正是此意。]

最近人们发现，4d $\cN=2$ SCFT 的涌现并不局限于 $\cN=2$ 超对称的 RG 流；它也可以来自某些
在中间能标只保持 $\cN=1$ 超对称、而在 IR 经历非平凡 SUSY 增强的流 [133,134]。这一情形中的
$\cN=1$ 规范理论通常与某个 $\cN=2$ SCFT 的幂零质量形变有关。

\SUB{2.6}{$\cN=2$ SCFT 的反常}

我们回顾：一个具有经典对称性 $G$（例如拉格朗日量的对称性）的量子场论，无论该对称性是规范
对称性还是整体对称性，都可能因量子效应而出现反常——这可以由路径积分测度的非不变性来探测。
关于 QFT 中反常基础的更多细节，读者可参阅综述 [135,136]。

规范反常意味着理论不自洽；而整体对称性的反常并不会使理论失效，恰恰相反，它是提取理论稳健
信息的强有力工具，尤其当理论是强耦合、且像许多 4d $\cN=2$ SCFT 那样没有拉格朗日量描述时。
我们把这类反常称为 't Hooft 反常：正如 't Hooft 所论证的 [137]，它沿 RG 流不变，因此在 UV
理论具有某些弱耦合描述时，可以在 UV 可靠地算出它。
% [边界] 本句（PDF p.27 起于「We refer to such an anomaly ... weakly」、PDF p.28 续完）已写完；
% 两页均属 tr-02 的块，p.28 起的正文见 tr02c.tex 开头。

% [Tutor Note: 「't Hooft 反常沿 RG 流不变」是本段后续所有论证的基础：它允许在 UV（弱耦合、有拉格朗日量）算出反常，再把它当作 IR 强耦合 SCFT 的精确数据使用（反常匹配）。注意这与「共形反常（Weyl 反常）的 $a,c$」不同——后者在 RG 流上会变（有 $a$-定理），所以第 2.6 节里两者必须分开处理：'t Hooft 反常给 UV/IR 匹配，$a,c$ 则要由 CB EFT 的数据重新拼出，见 (2.50)。]
